\documentclass[11pt]{article}

\usepackage[a4paper,margin=30mm]{geometry}
\usepackage[T1]{fontenc}
\usepackage[utf8]{inputenc}
\usepackage{lmodern}
\usepackage{microtype}
\usepackage{amsmath,amssymb,amsthm,mathtools}
\usepackage{booktabs,tabularx,array}
\usepackage{enumitem}
\usepackage{xcolor}
\usepackage[round,authoryear]{natbib}
\usepackage[colorlinks=true,linkcolor=blue!55!black,citecolor=red!45!black,urlcolor=blue!55!black]{hyperref}
\usepackage[nameinlink,noabbrev]{cleveref}

\hypersetup{
  pdftitle={A Soft Procrustes Formula for Gaussian KL-Unbalanced Optimal Transport},
  pdfauthor={Gabriel Peyr\'e},
  pdfsubject={Gaussian KL-unbalanced optimal transport, Bures geometry, and soft Procrustes alignment},
  pdfkeywords={optimal transport, unbalanced optimal transport, Gaussian measures, Bures metric, Procrustes, LogDet divergence}
}

\definecolor{deepblue}{RGB}{34,72,112}
\definecolor{warmgray}{RGB}{245,243,239}
\definecolor{darkred}{RGB}{132,55,45}

\newtheorem{theorem}{Theorem}[section]
\newtheorem{proposition}[theorem]{Proposition}
\newtheorem{lemma}[theorem]{Lemma}
\newtheorem{corollary}[theorem]{Corollary}
\theoremstyle{definition}
\newtheorem{definition}[theorem]{Definition}
\newtheorem{example}[theorem]{Example}
\theoremstyle{remark}
\newtheorem{remark}[theorem]{Remark}

\newcommand{\R}{\mathbb R}
\newcommand{\Spp}{\mathbb S_{++}^d}
\newcommand{\Spsd}{\mathbb S_{+}^d}
\newcommand{\Mplus}{\mathcal M_+}
\newcommand{\Ptwo}{\mathcal P_2}
\newcommand{\Id}{\mathrm{Id}}
\newcommand{\tr}{\operatorname{tr}}
\newcommand{\KL}{\operatorname{KL}}
\newcommand{\UOT}{\mathsf U}
\newcommand{\UB}{\mathsf{UB}}
\newcommand{\BW}{\mathsf B}
\newcommand{\cN}{\mathcal N}
\newcommand{\dd}{\,\mathrm d}
\newcommand{\norm}[1]{\left\lVert #1\right\rVert}

\title{\textbf{A Soft Procrustes Formula for Gaussian\\
KL-Unbalanced Optimal Transport}\\
\large Relaxed covariance factors, Riccati equations, and the Bures limit}
\author{Gabriel Peyr\'e}
\date{July 2026}

\begin{document}
\maketitle

\begin{abstract}
Quadratic Wasserstein transport between Gaussian measures decomposes into a
Euclidean cost between means and the Bures cost between covariances. The latter
is exactly an orthogonal Procrustes problem between matrix square roots. This
note studies the corresponding question for quadratic unbalanced optimal
transport with Kullback--Leibler penalties on the two marginals and no entropy
penalty on the coupling. After separating transported mass, the Gaussian
problem is an infimal convolution of Wasserstein distance with two Gaussian KL
divergences. We show that its covariance term has an exact \emph{soft
Procrustes} formulation: the hard Gram constraints of Bures are relaxed by
LogDet penalties, or equivalently the orthogonal factors are replaced by
general linear factors penalized for departing from the orthogonal group. A
single orthogonal Procrustes alignment remains after fixing two positive
stretches, but finite KL strength does not reduce to one hard orthogonal
alignment. We connect this representation to the known Riccati closed form,
treat unequal marginal penalties, recover ordinary Bures--Procrustes as the
large-penalty limit, and clarify which convexity properties survive.
\end{abstract}

\section{Question and answer}

The balanced Gaussian formula suggests asking whether the KL-relaxed problem
also comes from aligning covariance factors. The answer is positive, with one
important qualification.

\medskip
\noindent
\colorbox{warmgray}{\parbox{0.95\linewidth}{
\textbf{Executive answer.}
Gaussian KL-unbalanced quadratic transport admits an exact Procrustes-type
formula, but it is a \emph{soft} rather than a hard Procrustes problem. If
$\Sigma_i^{1/2}$ are prescribed covariance roots, the factors
$\Sigma_i^{1/2}X_i$ are free and the departure of $X_i$ from orthogonality is
penalized by
\[
  \Psi(X_i)=\norm{X_i}_{\mathrm F}^2
  -\log\det(X_iX_i^\top)-d.
\]
This penalty is non-negative and vanishes exactly on the orthogonal group.
For fixed positive stretches, the remaining rotation is the usual one-SVD
Procrustes alignment. In general, however, the stretches must also be
optimized, so there is no reduction to a single hard orthogonal alignment.
The optimizer is nevertheless explicit through the same Riccati equation as
the Gaussian KL-UOT closed form.}}
\medskip

The strictly entropic problem, which adds
$\epsilon\KL(\pi\mid\alpha\otimes\beta)$ to the coupling objective, is not the
object considered here. Its Gaussian solution is also explicit
\citep{JanatiEtAl2020}, but the optimal joint Gaussian is non-degenerate. Here
there is no coupling entropy: the optimizer is a scaled Wasserstein graph
coupling between two adjusted Gaussian marginals. The symmetric endpoint can
be recovered from the Gaussian formulas of \citet{JanatiEtAl2020}; the two
independent penalties are treated directly by \citet{YangZhang2026}.

\section{KL-unbalanced quadratic transport}

For finite non-negative measures $\rho$ and $\eta$, use the generalized
relative entropy
\begin{equation}\label{eq:generalized-kl}
 \KL(\rho\mid\eta)
 =\begin{cases}
   \displaystyle\int(r\log r-r+1)\dd\eta,
      & \rho=r\eta,\\[1mm]
   +\infty,&\rho\not\ll\eta.
  \end{cases}
\end{equation}
For a common marginal penalty $\tau>0$, define
\begin{equation}\label{eq:kl-uot}
 \UOT_\tau(\alpha_0,\alpha_1)
 =\inf_{\pi\in\Mplus(\R^d\times\R^d)}
 \left\{
   \int\norm{x-y}^2\dd\pi(x,y)
   +\tau\KL(\pi_0\mid\alpha_0)
   +\tau\KL(\pi_1\mid\alpha_1)
 \right\}.
\end{equation}
With the normalization in \eqref{eq:kl-uot},
$\UOT_\tau^{1/2}$ is the Gaussian Hellinger--Kantorovich endpoint metric on
finite measures over Euclidean space. It is related to, but should not be
confused with, the intrinsic Wasserstein--Fisher--Rao length metric; see
\citet{LieroMielkeSavare2018}. This note only studies the static cost
\eqref{eq:kl-uot}; see also \citet{ChizatEtAl2018} for the broader unbalanced
transport framework.

\subsection{Mass--shape separation}

Let $\alpha_i=a_i\mu_i$, where $a_i>0$ and $\mu_i$ are probability measures.
Introduce the normalized shape envelope
\begin{equation}\label{eq:shape-envelope}
 \mathcal E_\tau(\mu_0,\mu_1)
 =\inf_{\rho_0,\rho_1\in\Ptwo(\R^d)}
 \left\{
   W_2^2(\rho_0,\rho_1)
   +\tau\KL(\rho_0\mid\mu_0)
   +\tau\KL(\rho_1\mid\mu_1)
 \right\}.
\end{equation}

\begin{proposition}[Mass--shape decomposition]\label{prop:mass-shape}
For probability measures $\mu_0,\mu_1$ for which
$\mathcal E_\tau(\mu_0,\mu_1)<\infty$,
\begin{equation}\label{eq:mass-shape-value}
 \UOT_\tau(a_0\mu_0,a_1\mu_1)
 =\tau\left[
   a_0+a_1-2\sqrt{a_0a_1}
   \exp\!\left(-\frac{\mathcal E_\tau(\mu_0,\mu_1)}{2\tau}\right)
 \right].
\end{equation}
If $(\rho_0^*,\rho_1^*)$ minimizes \eqref{eq:shape-envelope}, the optimal
transported mass is
\begin{equation}\label{eq:optimal-mass}
 M^*=\sqrt{a_0a_1}
 \exp\!\left(-\frac{\mathcal E_\tau(\mu_0,\mu_1)}{2\tau}\right).
\end{equation}
The optimal plan is $M^*$ times a $W_2$-optimal probability coupling between
$\rho_0^*$ and $\rho_1^*$.
\end{proposition}

\begin{proof}
Write a nonzero plan as $\pi=M\bar\pi$, where $M$ is its total mass and
$\bar\pi$ is a probability coupling with marginals $\rho_0,\rho_1$. The
homogeneity identity
\[
 \KL(M\rho_i\mid a_i\mu_i)
 =M\KL(\rho_i\mid\mu_i)+M\log(M/a_i)-M+a_i
\]
reduces the objective, after minimizing $\bar\pi$ at fixed marginals, to
\begin{align*}
 M\mathcal A(\rho_0,\rho_1)
 &+\tau\{M\log(M/a_0)-M+a_0\}\\
 &+\tau\{M\log(M/a_1)-M+a_1\},
\end{align*}
where $\mathcal A$ is the expression minimized in
\eqref{eq:shape-envelope}. Differentiation in $M$ gives
$M=\sqrt{a_0a_1}e^{-\mathcal A/(2\tau)}$. Substitution gives
$\tau(a_0+a_1-2M)$, which is increasing in $\mathcal A$. Minimizing the
remaining shape variables proves both claims.
\end{proof}

The nonlinear exponential in \eqref{eq:mass-shape-value} is essential. The
shape envelope $\mathcal E_\tau$ is not itself the squared unbalanced
distance; it is the energy whose exponential controls how much mass is worth
transporting.

\section{Exact Gaussian reduction}

Assume from now on that
\begin{equation}\label{eq:gaussian-inputs}
 \mu_i=\cN(m_i,\Sigma_i),
 \qquad m_i\in\R^d,\quad \Sigma_i\in\Spp.
\end{equation}
For $P,\Sigma\in\Spp$, define the LogDet divergence
\begin{equation}\label{eq:logdet-divergence}
 D_{\mathrm{ld}}(P\mid\Sigma)
 =\tr(\Sigma^{-1}P)-\log\det(\Sigma^{-1}P)-d.
\end{equation}
The Gaussian KL formula reads
\begin{equation}\label{eq:gaussian-kl}
 \KL\bigl(\cN(u,P)\mid\cN(m,\Sigma)\bigr)
 =\frac12\left(
   \norm{u-m}_{\Sigma^{-1}}^2+D_{\mathrm{ld}}(P\mid\Sigma)
 \right).
\end{equation}

\begin{proposition}[Gaussian reduction]\label{prop:gaussian-reduction}
The infimum in \eqref{eq:shape-envelope} is attained by Gaussian measures.
More precisely,
\begin{align}
 \mathcal E_\tau(\mu_0,\mu_1)
 =\min_{\substack{u_0,u_1\in\R^d\\P_0,P_1\in\Spp}}
 \biggl\{&
   \norm{u_0-u_1}^2+\BW^2(P_0,P_1)\notag\\
 &+\frac\tau2\left(
   \norm{u_0-m_0}_{\Sigma_0^{-1}}^2
   +D_{\mathrm{ld}}(P_0\mid\Sigma_0)
 \right)\notag\\
 &+\frac\tau2\left(
   \norm{u_1-m_1}_{\Sigma_1^{-1}}^2
   +D_{\mathrm{ld}}(P_1\mid\Sigma_1)
 \right)\biggr\},
 \label{eq:finite-gaussian-envelope}
\end{align}
where
\begin{equation}\label{eq:bures-definition}
 \BW^2(P_0,P_1)
 =\tr P_0+\tr P_1
 -2\tr\!\left(P_0^{1/2}P_1P_0^{1/2}\right)^{1/2}.
\end{equation}
\end{proposition}

\begin{proof}
Let $\rho_i$ have mean $u_i$ and covariance $P_i$. Gelbrich's inequality
gives
\[
 W_2^2(\rho_0,\rho_1)
 \geq W_2^2\bigl(\cN(u_0,P_0),\cN(u_1,P_1)\bigr).
\]
Moreover, because the logarithm of a ratio of Gaussian densities is quadratic,
moment matching gives the Pythagorean identity
\[
 \KL(\rho_i\mid\mu_i)
 =\KL\bigl(\rho_i\mid\cN(u_i,P_i)\bigr)
 +\KL\bigl(\cN(u_i,P_i)\mid\mu_i\bigr).
\]
Replacing each $\rho_i$ by its moment-matching Gaussian can therefore only
decrease the objective. The Gaussian $W_2$ and KL formulas then give
\eqref{eq:finite-gaussian-envelope}. This argument is the Gaussian reduction
used in \citet{YangZhang2026}; see also \citet{Gelbrich1990}.
\end{proof}

The mean and covariance variables in \eqref{eq:finite-gaussian-envelope}
separate. Writing $\Delta m=m_0-m_1$, direct quadratic minimization gives
\begin{equation}\label{eq:mean-envelope-value}
 \mathcal E_{\tau,\mathrm{mean}}
 =\frac\tau2\,
 \Delta m^\top
 \left(\Sigma_0+\Sigma_1+\frac\tau2\Id\right)^{-1}
 \Delta m.
\end{equation}
The nontrivial part is therefore the covariance envelope
\begin{equation}\label{eq:covariance-envelope}
 \mathcal C_\tau(\Sigma_0,\Sigma_1)
 =\min_{P_0,P_1\in\Spp}
 \left\{
   \BW^2(P_0,P_1)
   +\frac\tau2D_{\mathrm{ld}}(P_0\mid\Sigma_0)
   +\frac\tau2D_{\mathrm{ld}}(P_1\mid\Sigma_1)
 \right\}.
\end{equation}

\section{The soft Procrustes representation}

Classical Bures geometry permits two equivalent factor formulations
\citep{BhatiaJainLim2019,MasarottoPanaretosZemel2019}. For $d'\geq d$,
\begin{equation}\label{eq:bures-free-factors}
 \BW^2(P_0,P_1)
 =\min_{\substack{F_i\in\R^{d\times d'}\\F_iF_i^\top=P_i}}
 \norm{F_0-F_1}_{\mathrm F}^2.
\end{equation}
If fixed factors $A_iA_i^\top=P_i$ are used instead, this is the orthogonal
Procrustes formula
\begin{equation}\label{eq:bures-hard-procrustes}
 \BW^2(P_0,P_1)
 =\min_{R\in\mathrm O(d')}\norm{A_0-A_1R}_{\mathrm F}^2.
\end{equation}
The KL-relaxed problem preserves \eqref{eq:bures-free-factors}, but replaces
the hard Gram constraints by penalties.

For a full-row-rank matrix $X\in\R^{d\times d'}$, define
\begin{equation}\label{eq:soft-stiefel-penalty}
 \Psi(X)
 =\norm{X}_{\mathrm F}^2-\log\det(XX^\top)-d.
\end{equation}

\begin{theorem}[Soft Procrustes formula]\label{thm:soft-procrustes}
Let $S_i=\Sigma_i^{1/2}$ and $d'\geq d$. Then
\begin{align}
 \mathcal C_\tau(\Sigma_0,\Sigma_1)
 =\min_{\substack{X_i\in\R^{d\times d'}\\
                   \operatorname{rank}X_i=d}}
 \biggl\{&
   \norm{S_0X_0-S_1X_1}_{\mathrm F}^2\notag\\
 &+\frac\tau2\Psi(X_0)+\frac\tau2\Psi(X_1)
 \biggr\}.
 \label{eq:soft-procrustes}
\end{align}
The penalty $\Psi$ is non-negative and
\begin{equation}\label{eq:soft-stiefel-zero-set}
 \Psi(X)=0\quad\Longleftrightarrow\quad XX^\top=\Id.
\end{equation}
Thus the hard row-Stiefel constraints of Bures--Procrustes are softened by
LogDet penalties. For $d'=d$, their zero set is exactly $\mathrm O(d)$.
\end{theorem}

\begin{proof}
Insert the free-factor identity \eqref{eq:bures-free-factors} into
\eqref{eq:covariance-envelope} and minimize jointly over $F_0,F_1$. Every
full-row-rank factor can be written $F_i=S_iX_i$. Cyclicity of the trace and
similarity invariance of the determinant give
\begin{align*}
 D_{\mathrm{ld}}(F_iF_i^\top\mid\Sigma_i)
 &=\tr(X_iX_i^\top)-\log\det(X_iX_i^\top)-d\\
 &=\Psi(X_i),
\end{align*}
which proves \eqref{eq:soft-procrustes}. If $s_1,\ldots,s_d$ are the singular
values of $X$, then
\[
 \Psi(X)=\sum_{j=1}^d(s_j^2-2\log s_j-1).
\]
Each summand is non-negative and vanishes only at $s_j=1$, proving
\eqref{eq:soft-stiefel-zero-set}.
\end{proof}

This identity can also be read directly in the original factor variables:
\begin{equation}\label{eq:penalized-gram-procrustes}
 \mathcal C_\tau
 =\min_{F_0,F_1}
 \left\{
   \norm{F_0-F_1}_{\mathrm F}^2
   +\frac\tau2D_{\mathrm{ld}}(F_0F_0^\top\mid\Sigma_0)
   +\frac\tau2D_{\mathrm{ld}}(F_1F_1^\top\mid\Sigma_1)
 \right\}.
\end{equation}
Ordinary Bures imposes $F_iF_i^\top=\Sigma_i$ exactly; KL-unbalanced Bures
penalizes violation of these two Gram constraints.

\subsection{One classical Procrustes step remains}

The soft formula is not a single orthogonal alignment, but it contains one.
Set $d'=d$ and use the left polar decompositions $X_i=H_iU_i$, with
$H_i\succ0$ and $U_i\in\mathrm O(d)$. The objective is invariant under
simultaneous right multiplication of $(X_0,X_1)$ by an orthogonal matrix.
Gauging out $U_0$ gives the following equivalent form.

\begin{corollary}[Stretched Procrustes formula]\label{cor:stretched-procrustes}
Define
\begin{equation}\label{eq:psi-spd}
 \psi(H)=\tr(H^2)-2\log\det H-d,
 \qquad H\in\Spp.
\end{equation}
Then
\begin{align}
 \mathcal C_\tau(\Sigma_0,\Sigma_1)
 =\min_{\substack{H_0,H_1\in\Spp\\R\in\mathrm O(d)}}
 \biggl\{&
   \norm{S_0H_0-S_1H_1R}_{\mathrm F}^2\notag\\
 &+\frac\tau2\psi(H_0)+\frac\tau2\psi(H_1)
 \biggr\}.
 \label{eq:stretched-procrustes}
\end{align}
For fixed $H_0,H_1$, the minimizer in $R$ is the ordinary orthogonal
Procrustes solution. Consequently,
\begin{align}
 \mathcal C_\tau
 =\min_{H_0,H_1\in\Spp}\biggl\{&
   \tr(\Sigma_0H_0^2)+\tr(\Sigma_1H_1^2)
   -2\norm{(S_1H_1)^\top(S_0H_0)}_*
   \notag\\
 &+\frac\tau2\psi(H_0)+\frac\tau2\psi(H_1)
 \biggr\},
 \label{eq:stretched-procrustes-nuclear}
\end{align}
where $\norm{\cdot}_*$ is the nuclear norm.
\end{corollary}

\begin{proof}
Right-multiply both polar factors by $U_0^\top$. This leaves all Frobenius
norms and both penalties unchanged, while replacing $(X_0,X_1)$ by
$(H_0,H_1R)$ with $R=U_1U_0^\top$. The first identity follows. For fixed
stretches, expanding the squared norm reduces the problem to maximizing a
trace over $R\in\mathrm O(d)$. The von Neumann trace inequality gives the
nuclear norm in \eqref{eq:stretched-procrustes-nuclear}; the maximizing
rotation is obtained from one singular-value decomposition
\citep{Schonemann1966}.
\end{proof}

The two positive stretches are the essential new degrees of freedom. Since
$\psi$ is not quadratic, the orbit argument behind
\eqref{eq:bures-hard-procrustes} does not eliminate them. Thus
\eqref{eq:stretched-procrustes} is an exact generalized Procrustes formula,
but not a one-SVD endpoint formula.

\subsection{Means, covariances, and masses in one formula}

The entire Gaussian problem admits the same interpretation. For
$\xi_i\in\R^d$ and $X_i\in\mathrm{GL}(d)$, define
\begin{align}
 \mathcal J_\tau(\xi_0,\xi_1,X_0,X_1)
 ={}&\norm{\Delta m+S_0\xi_0-S_1\xi_1}^2
 +\norm{S_0X_0-S_1X_1}_{\mathrm F}^2\notag\\
 &+\frac\tau2\left(
   \norm{\xi_0}^2+\norm{\xi_1}^2
   +\Psi(X_0)+\Psi(X_1)
 \right).
 \label{eq:full-soft-procrustes-energy}
\end{align}

\begin{corollary}[Gaussian KL-UOT as exponential soft Procrustes]
\label{cor:full-soft-procrustes}
For $\alpha_i=a_i\cN(m_i,\Sigma_i)$,
\begin{align}
 \UOT_\tau(\alpha_0,\alpha_1)
 =\min_{\xi_0,\xi_1,X_0,X_1}
 \tau\biggl[&a_0+a_1
 -2\sqrt{a_0a_1}\notag\\[-1mm]
 &\times\exp\!\left(
   -\frac{\mathcal J_\tau(\xi_0,\xi_1,X_0,X_1)}{2\tau}
 \right)\biggr].
 \label{eq:full-soft-procrustes-uot}
\end{align}
The minimization is over $\xi_i\in\R^d$ and $X_i\in\mathrm{GL}(d)$.
\end{corollary}

\begin{proof}
Write $u_i=m_i+S_i\xi_i$ in
\eqref{eq:finite-gaussian-envelope}, use
\cref{thm:soft-procrustes}, and then apply
\cref{prop:mass-shape}. The scalar function
$s\mapsto\tau(a_0+a_1-2\sqrt{a_0a_1}e^{-s/(2\tau)})$ is increasing, so the
minimum can be placed outside as in \eqref{eq:full-soft-procrustes-uot}.
\end{proof}

For centered unit-mass inputs, the restriction
\begin{equation}\label{eq:unbalanced-bures-definition}
 \UB_\tau^2(\Sigma_0,\Sigma_1)
 =\UOT_\tau\bigl(\cN(0,\Sigma_0),\cN(0,\Sigma_1)\bigr)
 =2\tau\left(1-e^{-\mathcal C_\tau/(2\tau)}\right)
\end{equation}
is the natural KL-unbalanced or Gaussian-Hellinger generalization of the
Bures metric. Formula \eqref{eq:soft-procrustes} is the requested
Procrustes-type representation, with the nonlinear conversion
\eqref{eq:unbalanced-bures-definition} from normalized shape energy to the
actual squared metric.

\section{Closed-form optimizer and the hidden Bures step}

The soft Procrustes problem is not merely a conceptual reformulation. Its
optimizer is the adjusted covariance pair appearing in the Gaussian KL-UOT
closed form. This section records the equal-penalty formula; the asymmetric
extension is given in \cref{sec:asymmetric}.

Set
\begin{equation}\label{eq:shifted-precisions}
 \eta=\frac2\tau,
 \qquad A_i=\Sigma_i^{-1},
 \qquad C_i=A_i+\eta\Id.
\end{equation}
Define
\begin{equation}\label{eq:closed-form-L}
 L_*=C_1^{-1/2}
 \left(C_1^{1/2}C_0C_1^{1/2}\right)^{1/2}
 C_1^{-1/2},
\end{equation}
and
\begin{equation}\label{eq:closed-form-PQ}
 P_*=[A_0+\eta(\Id-L_*)]^{-1},
 \qquad Q_*=L_*P_*L_*.
\end{equation}
For the means, let
\begin{equation}\label{eq:closed-form-means}
 h_*=[\Id+\eta(\Sigma_0+\Sigma_1)]^{-1}\Delta m,
 \quad
 u_0^*=m_0-\eta\Sigma_0h_*,
 \quad
 u_1^*=m_1+\eta\Sigma_1h_*.
\end{equation}

\begin{proposition}[Riccati solution of soft Procrustes]
\label{prop:riccati-soft-procrustes}
The matrices $P_*,Q_*$ in \eqref{eq:closed-form-PQ} are positive definite and
uniquely minimize \eqref{eq:covariance-envelope}. The normalized shape value is
\begin{align}
 \mathcal E_\tau(\mu_0,\mu_1)
 ={}&\Delta m^\top[\Id+\eta(\Sigma_0+\Sigma_1)]^{-1}\Delta m\notag\\
 &+\frac\tau2\log
 \frac{\det\Sigma_0\det\Sigma_1}{\det P_*\det Q_*}.
 \label{eq:closed-form-shape-value}
\end{align}
One minimizing factor pair in \eqref{eq:soft-procrustes} is
\begin{equation}\label{eq:closed-form-soft-factors}
 X_0^*=S_0^{-1}P_*^{1/2},
 \qquad
 X_1^*=S_1^{-1}L_*P_*^{1/2}.
\end{equation}
\end{proposition}

\begin{proof}
Let $L\succ0$ be the Gaussian optimal map from $P$ to $Q$, so $Q=LPL$.
The differential of $\BW^2(P,Q)$ with respect to $P$ is $\Id-L$, and its
differential with respect to $Q$ is $\Id-L^{-1}$. The first-order equations
for \eqref{eq:covariance-envelope} are therefore
\begin{equation}\label{eq:covariance-stationarity-note}
 P^{-1}=A_0+\eta(\Id-L),
 \qquad
 Q^{-1}=A_1+\eta(\Id-L^{-1}).
\end{equation}
Using $Q^{-1}=L^{-1}P^{-1}L^{-1}$ gives the Riccati equation
\begin{equation}\label{eq:riccati-equal}
 LC_1L=C_0.
\end{equation}
Its unique positive-definite solution is \eqref{eq:closed-form-L}, and back
substitution gives \eqref{eq:closed-form-PQ}. Global optimality and uniqueness
follow either from the KL-UOT dual certificate in \citet{YangZhang2026}, or
from joint convexity of the Bures term and strict convexity of the two LogDet
terms; see \cref{prop:convex-inner-problem} below.

At stationarity, taking traces in \eqref{eq:covariance-stationarity-note}
cancels the trace part of both Gaussian KL divergences against the Bures term.
Only the displayed log-determinants remain. Combining them with the quadratic
mean minimum gives \eqref{eq:closed-form-shape-value}. Finally, setting
$F_0=P_*^{1/2}$ and $F_1=L_*P_*^{1/2}$ gives
$F_iF_i^\top=P_i^*$ and
\[
 \norm{F_0-F_1}_{\mathrm F}^2=\BW^2(P_*,Q_*).
\]
Writing $F_i=S_iX_i^*$ yields \eqref{eq:closed-form-soft-factors}.
\end{proof}

Equation \eqref{eq:riccati-equal} reveals a second, internal Procrustes
structure: $L_*$ is exactly the Gaussian optimal map from the shifted precision
matrix $C_1$ to $C_0$. Hence the finite-penalty solution does use an ordinary
Bures square root or Procrustes alignment, but it applies to shifted
precisions and must be followed by the back-substitution
\eqref{eq:closed-form-PQ}.

\section{What survives from ordinary Bures geometry?}

\subsection{Large marginal penalties}

The soft formulation makes the balanced limit transparent. Since
$\Psi(X)\geq0$ with equality exactly on the row-Stiefel manifold, increasing
$\tau$ enforces $X_iX_i^\top\to\Id$ and $\xi_i\to0$. In the square case this
recovers the hard orthogonal factors of Bures--Procrustes.

\begin{proposition}[Bures--Procrustes limit]\label{prop:bures-limit}
For Gaussian probability measures $\mu_i=\cN(m_i,\Sigma_i)$,
\begin{equation}\label{eq:shape-bures-limit}
 \mathcal E_\tau(\mu_0,\mu_1)
 \longrightarrow W_2^2(\mu_0,\mu_1)
 \qquad(\tau\to\infty).
\end{equation}
Consequently, if $a_0=a_1=a$,
\begin{equation}\label{eq:uot-bures-limit}
 \UOT_\tau(a\mu_0,a\mu_1)
 \longrightarrow aW_2^2(\mu_0,\mu_1).
\end{equation}
If the masses differ, the leading term is
$\tau(\sqrt{a_0}-\sqrt{a_1})^2$.
\end{proposition}

\begin{proof}
The upper bound in \eqref{eq:shape-bures-limit} follows by choosing
$\rho_i=\mu_i$ in \eqref{eq:shape-envelope}. Bounded shape energy forces the
two KL terms to vanish as $\tau\to\infty$, hence any limit of minimizers has
marginals $\mu_0,\mu_1$; lower semicontinuity gives the converse bound. The
same conclusion follows directly from \eqref{eq:full-soft-procrustes-energy}:
bounded energy forces $\xi_i\to0$ and $X_iX_i^\top\to\Id$, after which the
remaining minimization is ordinary Gaussian Bures--Procrustes. Expanding the
exponential in \eqref{eq:mass-shape-value} proves the last two claims.
\end{proof}

\subsection{Convexity: inner problem versus endpoint parameters}

The word ``Procrustes'' often suggests a nonconvex matrix factorization. The
natural covariance variables show that the adjusted-covariance problem is in
fact convex.

\begin{proposition}[Convexity of the adjusted-covariance problem]
\label{prop:convex-inner-problem}
For fixed $\Sigma_0,\Sigma_1\succ0$, the objective in
\eqref{eq:covariance-envelope} is strictly jointly convex in $(P_0,P_1)$.
It therefore has a unique minimizer. The factor objective
\eqref{eq:soft-procrustes} is a redundant, nonconvex parameterization of this
convex problem and is invariant under simultaneous right orthogonal actions.
\end{proposition}

\begin{proof}
The squared Bures function is jointly convex on pairs of positive
semidefinite matrices \citep{BhatiaJainLim2019}. For fixed $\Sigma$, the map
$P\mapsto D_{\mathrm{ld}}(P\mid\Sigma)$ is strictly convex because its only
nonlinear term is $-\log\det P$, whose Hessian is positive definite on
symmetric directions. The sum is therefore strictly jointly convex. The
factor invariance follows from
$(F_0,F_1)\mapsto(F_0R,F_1R)$ for $R\in\mathrm O(d')$.
\end{proof}

This statement should not be confused with joint convexity in the prescribed
covariances $(\Sigma_0,\Sigma_1)$. Unlike squared Bures distance, the Gaussian
restriction of KL-UOT is not jointly convex in these parameters, even in one
dimension.

\begin{example}[Failure of covariance joint convexity]\label{ex:nonconvex-input-cov}
Take $d=1$, centered unit-mass Gaussians and $\tau=1$. Write $s_i>0$ for the
variance. Set
\[
 L=\sqrt{\frac{s_0^{-1}+2}{s_1^{-1}+2}},
 \qquad
 P=[s_0^{-1}+2(1-L)]^{-1},
 \qquad Q=L^2P.
\]
Then \eqref{eq:unbalanced-bures-definition} becomes
\begin{equation}\label{eq:scalar-unbalanced-bures}
 \UB_1^2(s_0,s_1)
 =2\left[1-\left(\frac{PQ}{s_0s_1}\right)^{1/4}\right].
\end{equation}
Direct substitution gives
\begin{align*}
 \UB_1^2(0.1,0.02)&=0.027134127639,\\
 \UB_1^2(0.3,0.04)&=0.089114415591,\\
 \UB_1^2(0.2,0.03)&=0.060423206395.
\end{align*}
The last pair is the midpoint of the first two, but
\[
 0.060423206395>
 \tfrac12(0.027134127639+0.089114415591)
 =0.058124271615.
\]
Thus covariance joint convexity fails. There is no contradiction with joint
convexity of KL-UOT in the input measures: a mixture of two Gaussian measures
is generally not the Gaussian whose covariance is the arithmetic mean.
\end{example}

\subsection{Why this is not one hard Procrustes problem}

The balanced Bures formula optimizes over a compact gauge orbit while keeping
the endpoint Gram matrices fixed. At finite $\tau$, the optimal factors in
\eqref{eq:soft-procrustes} generally satisfy
$X_iX_i^\top\neq\Id$: their singular values encode the KL-optimal covariance
relaxation. The nonquadratic LogDet penalty prevents the same orbit argument
from removing these stretches. One may compute the best rotation by one SVD
for each pair $(H_0,H_1)$, but the stretches remain coupled through
\eqref{eq:stretched-procrustes-nuclear}.

This does not constitute a theorem excluding every conceivable
finite-dimensional lift. It establishes the sharper statement relevant to the
standard Bures derivation: the natural square-root lift is exact, but it is a
penalized $\mathrm{GL}(d)$ alignment rather than a quotient by a single
orthogonal orbit. We did not find this soft-Procrustes identity stated in the
Gaussian KL-UOT references; it follows directly from their Gaussian reduction
combined with the classical free-factor Bures formula.

\section{Unequal marginal penalties}\label{sec:asymmetric}

The same representation applies to the asymmetric problem
\begin{equation}\label{eq:asymmetric-uot}
 \UOT_{\tau_0,\tau_1}(\alpha_0,\alpha_1)
 =\inf_\pi\left\{
   \int\norm{x-y}^2\dd\pi
   +\tau_0\KL(\pi_0\mid\alpha_0)
   +\tau_1\KL(\pi_1\mid\alpha_1)
 \right\}.
\end{equation}
This is no longer symmetric when $\tau_0\neq\tau_1$, so it should be called a
cost rather than a distance. Put $T=\tau_0+\tau_1$ and
$\theta_i=\tau_i/T$. Define
\begin{align}
 \mathcal J_{\tau_0,\tau_1}
 ={}&\norm{\Delta m+S_0\xi_0-S_1\xi_1}^2
 +\norm{S_0X_0-S_1X_1}_{\mathrm F}^2\notag\\
 &+\frac{\tau_0}{2}\{\norm{\xi_0}^2+\Psi(X_0)\}
 +\frac{\tau_1}{2}\{\norm{\xi_1}^2+\Psi(X_1)\}.
 \label{eq:asymmetric-soft-energy}
\end{align}
The mass calculation in \cref{prop:mass-shape} gives the exact asymmetric
soft-Procrustes formula
\begin{align}
 \UOT_{\tau_0,\tau_1}(a_0\mu_0,a_1\mu_1)
 =\min_{\xi_i,X_i}\biggl\{&
 \tau_0a_0+\tau_1a_1\notag\\[-1mm]
 &-T a_0^{\theta_0}a_1^{\theta_1}
 \exp\!\left(-\frac{\mathcal J_{\tau_0,\tau_1}}{T}\right)
 \biggr\}.
 \label{eq:asymmetric-soft-uot}
\end{align}

For completeness, the closed form of \citet{YangZhang2026} can be recovered
from this representation. Set
\begin{equation}\label{eq:asymmetric-shifts}
 \eta_i=\frac2{\tau_i},
 \qquad A_i=\Sigma_i^{-1},
 \qquad C_i=A_i+\eta_i\Id,
 \qquad \kappa=\eta_1-\eta_0.
\end{equation}
Define
\begin{align}
 S_*&=\frac12\left[
   \kappa\Id+
   \left(\kappa^2\Id+4C_1^{1/2}C_0C_1^{1/2}\right)^{1/2}
 \right],\notag\\
 L_*&=C_1^{-1/2}S_*C_1^{-1/2},
 \label{eq:asymmetric-L}
\end{align}
and
\begin{equation}\label{eq:asymmetric-PQ}
 P_*=[A_0+\eta_0(\Id-L_*)]^{-1},
 \qquad Q_*=L_*P_*L_*.
\end{equation}
Then $L_*$ is the unique positive solution of
\begin{equation}\label{eq:asymmetric-riccati}
 L_*C_1L_*-\kappa L_*=C_0.
\end{equation}
For the means,
\begin{equation}\label{eq:asymmetric-means}
 h_*=[\Id+\eta_0\Sigma_0+\eta_1\Sigma_1]^{-1}\Delta m,
 \quad
 u_*=m_0-\eta_0\Sigma_0h_*,
 \quad
 v_*=m_1+\eta_1\Sigma_1h_*.
\end{equation}
The minimum normalized energy is
\begin{align}
 \mathcal E_*
 ={}&\Delta m^\top
 [\Id+\eta_0\Sigma_0+\eta_1\Sigma_1]^{-1}\Delta m\notag\\
 &+\frac{\tau_0}{2}\log\frac{\det\Sigma_0}{\det P_*}
 +\frac{\tau_1}{2}\log\frac{\det\Sigma_1}{\det Q_*}.
 \label{eq:asymmetric-energy-value}
\end{align}
The transported mass and final value are
\begin{equation}\label{eq:asymmetric-mass-value}
 M_*=a_0^{\theta_0}a_1^{\theta_1}e^{-\mathcal E_*/T},
 \qquad
 \UOT_{\tau_0,\tau_1}=\tau_0a_0+\tau_1a_1-TM_*.
\end{equation}
When $\tau_0=\tau_1$, $\kappa=0$ and
\eqref{eq:asymmetric-riccati} reduces to the internal Bures equation
\eqref{eq:riccati-equal}. Unequal penalties deform this equation by a linear
Riccati term.

\section{Conclusion}

The KL-unbalanced Gaussian endpoint cost does have a precise analogue of the
Bures--Procrustes formula. It is best understood through three nested layers:
\begin{enumerate}[leftmargin=2em]
 \item total transported mass is eliminated analytically and produces the
 exponential conversion \eqref{eq:mass-shape-value};
 \item the normalized Gaussian problem chooses two adjusted means and
 covariances by KL penalization; and
 \item the covariance choice is exactly a soft Procrustes alignment of free
 square-root factors, with LogDet penalties replacing hard Gram constraints.
\end{enumerate}
At finite penalty, the correct symmetry is therefore not a larger hard
orthogonal group. It is a simultaneous orthogonal gauge acting on two general
linear factors whose singular values are optimized. One ordinary Procrustes
step survives for fixed stretches, and the complete optimizer is available
through a positive-definite Riccati equation. In the large-penalty limit, the
stretches converge to orthogonal factors and the whole construction collapses
to the classical Bures formula.

\bibliographystyle{plainnat}
\bibliography{references}

\end{document}
